% Part: sets-functions-relations
% Chapter: sets
% Section: russells-paradox

\documentclass[../../../include/open-logic-section]{subfiles}


\begin{document}

\olfileid{sfr}{set}{rus}
\olsection{د رسل پاراډوکس}

د غړو له مخې برابري د هغو $x$ ګانو د \emph{هم هغه} سټ دپاره د $\Setabs{x}{\phi(x)}$ نښې د کارولو جواز راکوي چې~$\phi(x)$ وي۔ خو د غړو له مخې برابري \emph{په اصل کښې} يوازې د دې خبرې جواز راکوي: \emph{که} داسې يو سټ وي چې غړي ئې ټول او يوازې هغه څيزونه وي چې $\phi$ وي، \emph{نو} داسې سټ يوازې يو وي۔ په بل عبارت: د يو~$\phi$ له ټاکلو وروسته، سټ $\Setabs{x}{\phi(x)}$ يوازينے دے، \emph{که موجود وي}۔

خو دا شرط مهم دے!{} بنيادي خبره دا ده چې له هر خاصيت څخه د \emph{سټ جوړول} ممکن نۀ دي۔ يعنې ځينې خاصيتونه سټونه \emph{نۀ} تعريفوي۔ که ټولو تعريفولے، نو له ښکاره تناقضونو سره به مخ شوو۔ د دې تر ټولو مشهوره بېلګه د رسل پاراډوکس دے۔

سټونه د نورو سټونو !!{element}s کېدے شي؛ د بېلګې په توګه، د يو سټ~$A$ د فرعي سټونو سټ له سټونو جوړ وي۔ نو دا پوښتل يا څېړل معقول دي چې يو سټ د بل سټ !!a{element} دے که نۀ۔ آيا يو سټ د خپل ځان غړے کېدے شي؟ د سټ په تصور کښې ظاهراً هېڅ داسې خبره نشته چې دا ناممکنه کړي۔ د بېلګې په توګه، که \emph{ټول} سټونه د څيزونو يوه ټولګه جوړوي، يو څوک ښائي فکر وکړي چې دا ټول په يو سټ کښې راټولېدے شي، يعنې د ټولو سټونو په سټ کښې۔ او دا چې پخپله يو سټ دے، نو د ټولو سټونو د سټ !!a{element} به وي۔

د رسل پاراډوکس هغه وخت راولاړېږي چې دا خاصيت په پام کښې ونيسو چې يو سټ خپل ځان د !!a{element} په توګه ونۀ لري، يعنې \emph{د خپل ځان غړے نۀ وي}۔ که فرض کړو چې د ټولو هغو سټونو يو سټ شته چې خپل ځان د !!a{element} په توګه نۀ لري، نو څه به وشي؟ آيا
\[
R = \Setabs{x}{x \notin x}
\]
شته؟ څرګندېږي چې ثابتولے شو دا نشته۔

\begin{thm}[د رسل پاراډوکس]\ollabel{thm:russells-paradox}
	د $R = \Setabs{x}{x \notin x}$ په بڼه هېڅ سټ نشته۔
\end{thm}

\begin{proof}
که $R = \Setabs{x}{x \notin x}$ موجود وي، نو $R \in R$ هله او يوازې هله وي چې $R \notin R$ وي، او دا يو تناقض دے۔
\end{proof}

\begin{tagblock}{novice}
\begin{explain}
راځئ دا ثبوت لږ ورو وګورو۔ که $R$ موجود وي، دا پوښتنه معقوله ده چې $R \in R$ دے که نۀ۔ فرض کړئ چې په رښتيا $R \in R$۔ اوس، $R$~د هغو ټولو سټونو د سټ په توګه تعريف شوے ؤ چې د خپلو ځانونو !!{element}s نۀ دي۔ نو که $R \in R$ وي، بيا $R$ پخپله د $R$ تعريفوونکے خاصيت نۀ لري۔ خو يوازې هغه سټونه په~$R$ کښې دي چې دا خاصيت لري؛ ځکه $R$ د~$R$ !!a{element} نۀ شي کېدے، يعنې $R \notin R$۔ خو $R$ په يو وخت د~$R$ !!a{element} هم نۀ شي کېدے او ناغړے هم، نو له تناقض سره مخ يو۔

څرنګه چې د $R \in R$ فرض تناقض ته رسوي، نو $R \notin R$ لرو۔ خو دا هم تناقض ته رسوي!{} ځکه که $R \notin R$ وي، نو $R$ پخپله د $R$ تعريفوونکے خاصيت لري، او ځکه $R$ به د هغو ټولو نورو سټونو په شان د $R$ !!a{element} وي چې د خپل ځان غړي نۀ دي۔ او بيا هم، په يو وخت د~$R$ ناغړے او !!a{element} دواړه نۀ شي کېدے۔
\end{explain}
\end{tagblock}

\begin{digress}
د سټونو داسې نظريه څنګه جوړه کړو چې د رسل له پاراډوکس څخه ځان وساتي، يعنې له دې \emph{متناقضې} دعوې څخه ځان وساتي چې $R = \Setabs{x}{x \notin x}$ موجود دے؟ د دې دپاره داسې بديهي اصلونه ټاکل پکار دي چې دا بيان کړي چې سټونه کله موجود وي (او کله نۀ وي)، او ډېر دقيق شرطونه ورکړي۔

د سټونو هغه نظريه چې په دې باب کښې په لنډه توګه بيان شوې، دا کار نۀ کوي۔ دا \emph{په رښتيا ساده نظريه} ده۔ دا يوازې درته وائي چې سټونه د غړو له مخې د برابرۍ پابند دي، او که ځينې سټونه ولرئ، د هغو اتحاد، تقاطع او داسې نور جوړولے شئ۔ د سټونو نظريه له دې نه په زيات دقت هم جوړېدل ممکن دي۔ \oliflabeldef{cumul:::part}{دا دقيق بحث به د \olref[cumul][][]{part} برخې دپاره پرېږدو۔ اوس به په ساده طريقه مخ ته ځو، او په احتياط به هڅه کوو چې له تناقضونو ځان وساتو۔}{}
\end{digress}


\end{document}
